% figures/w02-windup.svg — 와인드업이 왜 적분기 잘못이 아닌가.
%
%   bash _tools/tikz2svg.sh figures/src/w02-windup.tex figures/w02-windup.svg
%
% (a) 포화하는 동안 루프가 끊긴다는 것을 블록선도로,
% (b) 요구가 한계 밖으로 얼마나 나갔는지를 포화 특성 위에 실척으로.
%
% 주의: TeX 제어어에는 숫자를 못 쓴다. \def\Y0{} 은 \Y 에 인자 0 을 붙인 것이
% 되어 조용히 틀린다 — 실제로 한 번 그랬다. 매크로 이름은 글자만 쓴다.
\documentclass[border=5pt]{standalone}
% gnc-style.tex — 이 볼트의 TikZ 그림이 공유하는 스타일.
%
% 저널 그림 규약을 스타일 하나로 굳혀 둔다. 그림마다 선 굵기나 화살촉을
% 다시 정하지 않는다 — 그것이 그림이 제각각으로 보이는 원인이다.
%
% 쓰는 법:  % gnc-style.tex — 이 볼트의 TikZ 그림이 공유하는 스타일.
%
% 저널 그림 규약을 스타일 하나로 굳혀 둔다. 그림마다 선 굵기나 화살촉을
% 다시 정하지 않는다 — 그것이 그림이 제각각으로 보이는 원인이다.
%
% 쓰는 법:  % gnc-style.tex — 이 볼트의 TikZ 그림이 공유하는 스타일.
%
% 저널 그림 규약을 스타일 하나로 굳혀 둔다. 그림마다 선 굵기나 화살촉을
% 다시 정하지 않는다 — 그것이 그림이 제각각으로 보이는 원인이다.
%
% 쓰는 법:  \input{gnc-style}  (figures/src/ 안에서)

\usepackage{tikz}
\usepackage{amsmath}
\usepackage{xcolor}
\usetikzlibrary{arrows.meta, positioning, calc, fit, backgrounds, decorations.pathmorphing}

% ---- 색 --------------------------------------------------------------------
% 색은 강조 하나에만 쓴다. 나머지는 검정.
\definecolor{ink}{HTML}{1A1A1A}
\definecolor{accent}{HTML}{7C3AED}
\definecolor{warn}{HTML}{B91C1C}
\definecolor{soft}{HTML}{666666}

% ---- 선과 화살촉 -----------------------------------------------------------
\tikzset{
  % 신호선. 화살촉은 작고 뾰족한 것 하나뿐이다.
  sig/.style   = {ink, line width=0.5pt, -{Stealth[length=4pt, width=3pt]}},
  wire/.style  = {ink, line width=0.5pt},
  % 강조 경로 — 파선으로 구분한다. 흑백 인쇄에서도 살아야 하기 때문이다.
  asig/.style  = {accent, line width=0.5pt, densely dashed,
                  -{Stealth[length=4pt, width=3pt]}},
  awire/.style = {accent, line width=0.5pt, densely dashed},
  % 블록. 흰 바탕, 직각, 검은 테두리.
  blk/.style   = {draw=ink, line width=0.55pt, fill=white,
                  minimum width=17mm, minimum height=8mm, inner sep=2pt, font=\small},
  % 합산점
  sum/.style   = {draw=ink, line width=0.55pt, fill=white, circle,
                  inner sep=0pt, minimum size=4.6mm, font=\scriptsize},
  asum/.style  = {draw=accent, line width=0.55pt, fill=white, circle,
                  inner sep=0pt, minimum size=4.6mm, font=\scriptsize},
  % 분기점
  dot/.style   = {ink, fill=ink, circle, inner sep=0pt, minimum size=1.5mm},
  % 신호 이름 — 선 위에 작게
  sname/.style = {font=\small, inner sep=1.5pt},
  % 그림 안의 짧은 설명. 그림 안의 글은 최소로 둔다
  note/.style  = {font=\scriptsize, soft, inner sep=1.5pt},
  anote/.style = {font=\scriptsize, accent, inner sep=1.5pt},
  % 축
  axis/.style  = {ink, line width=0.5pt},
  tick/.style  = {font=\scriptsize, soft},
}

% 캡션 한 줄 — 그림 아래 가는 선과 함께
\newcommand{\gnccaption}[2]{%
  \node[anchor=north west, text width=#1, align=left, font=\scriptsize, ink]
        at (current bounding box.south west) {\vspace{2pt}#2};}
  (figures/src/ 안에서)

\usepackage{tikz}
\usepackage{amsmath}
\usepackage{xcolor}
\usetikzlibrary{arrows.meta, positioning, calc, fit, backgrounds, decorations.pathmorphing}

% ---- 색 --------------------------------------------------------------------
% 색은 강조 하나에만 쓴다. 나머지는 검정.
\definecolor{ink}{HTML}{1A1A1A}
\definecolor{accent}{HTML}{7C3AED}
\definecolor{warn}{HTML}{B91C1C}
\definecolor{soft}{HTML}{666666}

% ---- 선과 화살촉 -----------------------------------------------------------
\tikzset{
  % 신호선. 화살촉은 작고 뾰족한 것 하나뿐이다.
  sig/.style   = {ink, line width=0.5pt, -{Stealth[length=4pt, width=3pt]}},
  wire/.style  = {ink, line width=0.5pt},
  % 강조 경로 — 파선으로 구분한다. 흑백 인쇄에서도 살아야 하기 때문이다.
  asig/.style  = {accent, line width=0.5pt, densely dashed,
                  -{Stealth[length=4pt, width=3pt]}},
  awire/.style = {accent, line width=0.5pt, densely dashed},
  % 블록. 흰 바탕, 직각, 검은 테두리.
  blk/.style   = {draw=ink, line width=0.55pt, fill=white,
                  minimum width=17mm, minimum height=8mm, inner sep=2pt, font=\small},
  % 합산점
  sum/.style   = {draw=ink, line width=0.55pt, fill=white, circle,
                  inner sep=0pt, minimum size=4.6mm, font=\scriptsize},
  asum/.style  = {draw=accent, line width=0.55pt, fill=white, circle,
                  inner sep=0pt, minimum size=4.6mm, font=\scriptsize},
  % 분기점
  dot/.style   = {ink, fill=ink, circle, inner sep=0pt, minimum size=1.5mm},
  % 신호 이름 — 선 위에 작게
  sname/.style = {font=\small, inner sep=1.5pt},
  % 그림 안의 짧은 설명. 그림 안의 글은 최소로 둔다
  note/.style  = {font=\scriptsize, soft, inner sep=1.5pt},
  anote/.style = {font=\scriptsize, accent, inner sep=1.5pt},
  % 축
  axis/.style  = {ink, line width=0.5pt},
  tick/.style  = {font=\scriptsize, soft},
}

% 캡션 한 줄 — 그림 아래 가는 선과 함께
\newcommand{\gnccaption}[2]{%
  \node[anchor=north west, text width=#1, align=left, font=\scriptsize, ink]
        at (current bounding box.south west) {\vspace{2pt}#2};}
  (figures/src/ 안에서)

\usepackage{tikz}
\usepackage{amsmath}
\usepackage{xcolor}
\usetikzlibrary{arrows.meta, positioning, calc, fit, backgrounds, decorations.pathmorphing}

% ---- 색 --------------------------------------------------------------------
% 색은 강조 하나에만 쓴다. 나머지는 검정.
\definecolor{ink}{HTML}{1A1A1A}
\definecolor{accent}{HTML}{7C3AED}
\definecolor{warn}{HTML}{B91C1C}
\definecolor{soft}{HTML}{666666}

% ---- 선과 화살촉 -----------------------------------------------------------
\tikzset{
  % 신호선. 화살촉은 작고 뾰족한 것 하나뿐이다.
  sig/.style   = {ink, line width=0.5pt, -{Stealth[length=4pt, width=3pt]}},
  wire/.style  = {ink, line width=0.5pt},
  % 강조 경로 — 파선으로 구분한다. 흑백 인쇄에서도 살아야 하기 때문이다.
  asig/.style  = {accent, line width=0.5pt, densely dashed,
                  -{Stealth[length=4pt, width=3pt]}},
  awire/.style = {accent, line width=0.5pt, densely dashed},
  % 블록. 흰 바탕, 직각, 검은 테두리.
  blk/.style   = {draw=ink, line width=0.55pt, fill=white,
                  minimum width=17mm, minimum height=8mm, inner sep=2pt, font=\small},
  % 합산점
  sum/.style   = {draw=ink, line width=0.55pt, fill=white, circle,
                  inner sep=0pt, minimum size=4.6mm, font=\scriptsize},
  asum/.style  = {draw=accent, line width=0.55pt, fill=white, circle,
                  inner sep=0pt, minimum size=4.6mm, font=\scriptsize},
  % 분기점
  dot/.style   = {ink, fill=ink, circle, inner sep=0pt, minimum size=1.5mm},
  % 신호 이름 — 선 위에 작게
  sname/.style = {font=\small, inner sep=1.5pt},
  % 그림 안의 짧은 설명. 그림 안의 글은 최소로 둔다
  note/.style  = {font=\scriptsize, soft, inner sep=1.5pt},
  anote/.style = {font=\scriptsize, accent, inner sep=1.5pt},
  % 축
  axis/.style  = {ink, line width=0.5pt},
  tick/.style  = {font=\scriptsize, soft},
}

% 캡션 한 줄 — 그림 아래 가는 선과 함께
\newcommand{\gnccaption}[2]{%
  \node[anchor=north west, text width=#1, align=left, font=\scriptsize, ink]
        at (current bounding box.south west) {\vspace{2pt}#2};}


\begin{document}
\begin{tikzpicture}[x=1mm, y=1mm]

% ===================== (a) the loop =========================================
\begin{scope}
  \node[font=\small\bfseries, anchor=west] at (-13, 17) {(a)\quad where the loop breaks};

  \node[sum]                                (s1)  at (0,0) {};
  \node[blk, right=8mm  of s1, yshift=8mm]  (kp)  {$K_p$};
  \node[blk, right=8mm  of s1, yshift=-8mm] (ki)  {$K_i/s$};
  \node[sum, right=40mm of s1]              (s2)  {};
  \node[blk, right=16mm of s2]              (sat) {};
  \node[blk, right=16mm of sat]             (pl)  {hull};

  % 포화 글리프 — 상자 안에 글자를 또 넣지 않는다
  \draw[wire] ([shift={(-5.5mm,-2.4mm)}]sat.center) -- ([shift={(-2mm,-2.4mm)}]sat.center)
              -- ([shift={( 2mm, 2.4mm)}]sat.center) -- ([shift={( 5.5mm, 2.4mm)}]sat.center);

  % 입력과 분기
  \draw[sig] (-13,0) -- node[sname, above, pos=0.25] {$u_d$} (s1);
  \coordinate (br) at ($(s1)+(4mm,0)$);
  \draw[wire] (s1) -- (br);
  \node[dot] at (br) {};
  \draw[wire] (br) |- ($(kp.west)+(-3mm,0)$);
  \draw[sig]  ($(kp.west)+(-3mm,0)$) -- (kp);
  \draw[wire] (br) |- ($(ki.west)+(-3mm,0)$);
  \draw[sig]  ($(ki.west)+(-3mm,0)$) -- (ki);
  \node[sname, above right=0mm and -0.5mm of br] {$e$};

  % 합류
  \coordinate (jn) at ($(s2)+(-5mm,0)$);
  \draw[wire] (kp) -| (jn);  \draw[wire] (ki) -| (jn);
  \draw[sig]  (jn) -- (s2);

  % 앞으로 가는 길
  \draw[sig] (s2)  -- node[sname, above] {$X_{\mathrm{cmd}}$} (sat);
  \draw[sig] (sat) -- node[sname, above, pos=0.28] {$X_{\mathrm{sat}}$} (pl);
  \coordinate (out) at ($(pl.east)+(9mm,0)$);
  \draw[wire] (pl) -- (out);
  \draw[sig]  (out) -- ++(7mm,0) node[sname, above, pos=0.45] {$u$};

  % 부호
  \node[font=\scriptsize] at ($(s1)+(-2.8mm, 3.0mm)$) {$+$};
  \node[font=\scriptsize] at ($(s1)+(-3.6mm,-3.2mm)$) {$-$};
  \node[font=\scriptsize] at ($(s2)+(-3.0mm, 3.0mm)$) {$+$};
  \node[font=\scriptsize] at ($(s2)+( 3.0mm, 3.0mm)$) {$+$};

  % 되먹임
  \node[dot] at (out) {};
  \draw[wire] (out) |- (0,-22);
  \draw[sig]  (0,-22) -- (s1);

  % 루프가 끊기는 자리 — 작게. 크게 그리면 그림이 포스터가 된다
  \coordinate (cut) at ($(sat.east)!0.80!(pl.west)$);
  \draw[warn, line width=0.7pt] ($(cut)+(-1.5mm,-1.5mm)$) -- ($(cut)+(1.5mm,1.5mm)$);
  \draw[warn, line width=0.7pt] ($(cut)+(-1.5mm, 1.5mm)$) -- ($(cut)+(1.5mm,-1.5mm)$);
  \node[font=\scriptsize, warn, below=0.8mm of cut] {cut};

  % 안티와인드업 — 진짜 구조대로: X_cmd 와 X_sat 의 차를 적분기로 되돌린다
  \coordinate (tc) at ($(s2)!0.42!(sat.west)$);
  \coordinate (ts) at ($(sat.east)!0.28!(pl.west)$);
  \node[dot] at (tc) {};  \node[dot] at (ts) {};
  \node[asum] (aw) at ($(tc)!0.5!(ts) + (0,-13mm)$) {};
  \draw[awire] (tc) |- ($(aw.west)+(-4mm,0)$);
  \draw[asig]  ($(aw.west)+(-4mm,0)$) -- (aw);
  \draw[awire] (ts) |- ($(aw.east)+(4mm,0)$);
  \draw[asig]  ($(aw.east)+(4mm,0)$) -- (aw);
  \node[font=\scriptsize, accent] at ($(aw)+(-3.0mm,3.0mm)$) {$+$};
  \node[font=\scriptsize, accent] at ($(aw)+( 3.0mm,3.0mm)$) {$-$};
  \draw[awire] (aw) -- ++(0,-4mm) -| ($(ki.south)+(0,-4mm)$);
  \draw[asig]  ($(ki.south)+(0,-4mm)$) -- (ki.south);
  \node[anote, right=2mm of aw] {$X_{\mathrm{cmd}}-X_{\mathrm{sat}}$};
\end{scope}

% ===================== (b) the saturation characteristic ====================
\begin{scope}[shift={(6,-70)}]
  \node[font=\small\bfseries, anchor=west] at (-19, 36)
       {(b)\quad how far outside the limit the demand went};

  % 데이터 -> mm.  x: 0..3700 N -> 0..64 mm ,  y: -220..320 N -> 0..30 mm
  \def\XS{0.01730}
  \def\YS{0.05556}
  \def\Yzero{12.22}          % y = 0 의 높이 = 220 * YS

  \draw[axis] (-3,\Yzero) -- (66,\Yzero);
  \draw[axis] (0,-1) -- (0,31);

  % 포화 특성
  \draw[ink, line width=0.8pt]
        (-3, {\Yzero + (-133.42)*\YS}) -- ({-133.42*\XS}, {\Yzero + (-133.42)*\YS})
     -- ({ 239.36*\XS}, {\Yzero + 239.36*\YS}) -- (64, {\Yzero + 239.36*\YS});
  \draw[soft, line width=0.35pt, densely dotted]
        (0, {\Yzero + 239.36*\YS}) -- (64, {\Yzero + 239.36*\YS});
  \node[note, anchor=west] at (16, {\Yzero + 239.36*\YS + 4.8})
       {$X_{\mathrm{sat}}\le 239.36$\,N \ \ (the flat part)};

  % 측정된 최대 요구
  \draw[accent, line width=0.5pt, densely dashed]
        ({3480*\XS}, \Yzero) -- ({3480*\XS}, {\Yzero + 239.36*\YS});
  \fill[accent] ({3480*\XS}, {\Yzero + 239.36*\YS}) circle (0.7mm);
  \node[anote, anchor=east] at ({3480*\XS - 2}, {\Yzero + 239.36*\YS - 4.0})
       {3480\,N asked for};
  \node[anote, anchor=east] at ({3480*\XS - 2}, {\Yzero + 239.36*\YS - 7.4})
       {239\,N delivered};

  % 눈금
  \foreach \v in {0, 1000, 2000, 3000} {
    \draw[axis] ({\v*\XS}, \Yzero) -- ({\v*\XS}, \Yzero-1.2);
    \node[tick, below=0.3mm] at ({\v*\XS}, \Yzero-1.2) {\v};
  }
  \node[font=\small, anchor=north] at (32, \Yzero-6) {$X_{\mathrm{cmd}}$\ \ [N]};
  % y 눈금 — 한계 두 개만. 축에 눈금이 없으면 "실척" 이라는 말이 근거를 잃는다
  \foreach \v/\lab in {0/0, 239.36/239, -133.42/{-133}} {
    \draw[axis] (0, {\Yzero + \v*\YS}) -- (-1.3, {\Yzero + \v*\YS});
    \node[tick, anchor=east] at (-1.8, {\Yzero + \v*\YS}) {\lab};
  }
  \node[font=\small, anchor=south, rotate=90] at (-11, 15) {$X_{\mathrm{sat}}$\ \ [N]};
\end{scope}

\end{tikzpicture}
\end{document}
